\section{Introduction}
\label{sec:introduction}

For a partition $\lambda\vdash n$, write $S^\lambda$ for the corresponding
irreducible complex representation of the symmetric group $S_n$. Its
Kronecker coefficients are
\[
 g(\lambda,\mu,\nu)
 =\dim\Hom_{S_n}(S^\nu,S^\lambda\otimes S^\mu),
\]
where the tensor product carries the diagonal action. We prove the
following statement.

\begin{theorem}\label{thm:main}
For every positive integer $n\notin\{2,4,9\}$ there is a partition
$\lambda\vdash n$ such that
\[
 g(\lambda,\lambda,\nu)>0\qquad\text{for every }\nu\vdash n.
\]
Thus one irreducible representation has a tensor square containing every
irreducible representation of $S_n$.
\end{theorem}

The partition $\lambda$ is fixed for each degree: it does not vary with
the desired constituent $\nu$. In particular the trivial and sign
representations both occur. The latter forces $\lambda$ to be
self-conjugate, since
\[
 g(\lambda,\lambda,(1^n))
 =\dim\Hom_{S_n}(S^\lambda,S^{\lambda^t})
 =\begin{cases}1,&\lambda=\lambda^t,\\0,&\lambda\ne\lambda^t.
 \end{cases}
\]
Our construction meets this necessary condition throughout.

\begin{corollary}[Unipotent tensor squares]\label{cor:unipotent-squares}
For each positive integer $n\notin\{2,4,9\}$, choose a partition
$\lambda_n$ supplied by Theorem~\ref{thm:main}, including its
finite-range witnesses. For every prime power $q$, let $U_q^\lambda$
denote the unipotent irreducible complex character of $\GL_n(\mathbb F_q)$
indexed by $\lambda\vdash n$. Then
$U_q^{\lambda_n}\otimes U_q^{\lambda_n}$ contains $U_q^\nu$
for every $\nu\vdash n$.
\end{corollary}

\begin{proof}
Theorem~\ref{thm:main} gives $g(\lambda_n,\lambda_n,\nu)>0$ for every
$\nu\vdash n$. Letellier's nonnegativity theorem and constant-term bound
\cite[Theorem~4 and Proposition~6]{Letellier}, recalled in
\cite[Theorem~2.8]{LetellierNam}, transfer this positivity to the
unipotent multiplicity for the same partition labels and every prime
power $q$.
\end{proof}

The Steinberg character already gives this unipotent coverage in every
degree \cite[Proposition~33]{Letellier}, and staircase labels give it
at triangular degrees \cite[Theorem~1.1]{LetellierNam}. The corollary
identifies further covering squares using the same labels that solve
the symmetric-group problem.

\subsection{History and significance}

Pak, Panova, and Vallejo formulate the assertion for $n\ge3$, with
$n\ne4,9$, as the \emph{tensor square conjecture} for symmetric groups
\cite[Conjecture~1.1]{PPV}. The case $n=1$ is immediate. Hence
Theorem~\ref{thm:main} resolves that conjecture affirmatively, including
every nontriangular degree.

The related \emph{Saxl conjecture} specifies the staircase partition
$\rho_m=(m,m-1,\ldots,1)$ at the triangular degree
$N_m=m(m+1)/2$; see \cite[Conjecture~1.2]{PPV}.
Saxl proposed it at the UCLA Combinatorics Seminar on 20 March 2012,
as recorded in \cite[Footnote~2]{PPV}.
It is a more specific assertion about a distinguished family, but only
at triangular degrees. Earlier work established substantial support
families: Pak, Panova, and Vallejo gave a character-nonvanishing test for
constituents of self-conjugate squares \cite[Lemma~1.3]{PPV},
while Ikenmeyer proved the staircase assertion for partitions comparable
with the staircase in dominance order \cite[Theorem~2.1]{Ikenmeyer}.
His triangular arrangement of alternating tensors and nonzero contractions
\cite[Sections~3--5]{Ikenmeyer} is a methodological predecessor of the
cyclic construction used here.
Heide, Saxl, Tiep, and Zalesski conjectured the analogous tensor-square
covering property for alternating simple groups \cite[Remark~1.3(3)]{HSTZ},
partly motivated by Thompson's conjugacy-class-square conjecture for
finite simple groups.

Several complementary approaches explain the strength of the square
problem. Luo and Sellke proved that staircase squares contain almost all
partitions under both the uniform and Plancherel measures, and obtained
full fourth-power coverage in all sufficiently large degrees, using
modified shapes at nontriangular degrees
\cite[Theorems~1.4--1.6]{LuoSellke}. Harman and Ryba subsequently proved
full coverage for every staircase tensor cube \cite[Theorem~1.1]{HarmanRyba}.
Spin-character methods gave Bessenrodt all double-hook constituents,
and Li proved the triple-hook case: these are targets with two and
three diagonal cells, respectively
\cite[Theorem~4.10]{Bessenrodt}\cite[Theorem~4.22]{Li}.
Modular methods yield, among other families, all
2-height-zero constituents and framed staircases
\cite[Theorems~5.2 and~5.5]{BessenrodtBowmanSutton}.
More recent semigroup and generalized-block criteria give additional
constituent families \cite[Theorem~1.6]{Ebrahimi}.
These results concern different support families or tensor powers:
neither almost-all coverage nor a third or fourth power gives the
universal square required here.

The passage to nontriangular degrees also has a direct antecedent.
Outside degrees $2,4,9$, Li conjectured universal-square support for every
self-conjugate staircase-like partition in a parity-dependent family,
and asked whether
suitable additions of one or two boxes preserve universal-square support
\cite[Section~5.5, Definition~5.3, Conjecture~1 and Problem~5.4]{Li}.
Our large-degree construction uses a different family, making matching
additions to the first two rows and columns of a parity-compatible
staircase. The proof controls all constituents of the chosen enlargement;
the candidate is fixed before the target is specified.
Bessenrodt and Bowman formulate stronger conjectural refinements in
terms of symmetric and alternating parts of tensor squares
\cite[Conjectures~10.2--10.3]{BessenrodtBowman}; the tensor product in
Theorem~\ref{thm:main} is the ordinary diagonal tensor square.

We use the stronger cyclic form of the staircase result proved in
\emph{A Cyclic Polytabloid Proof of Saxl's Conjecture}
\cite[Theorem~3.1]{SaxlCyclic}.
The exact theorem and generator are specified in
Theorem~\ref{thm:cyclic-input}. The new task here is to pass from staircase
degrees to all the degrees in Theorem~\ref{thm:main}.

\subsection{The proof and its reusable ingredients}

The obstacle is not simply to place a smaller staircase inside a larger
diagram. Its constituents must survive a map from the actual enlarged
square, and attachments can introduce cancelling alternating terms.
Section~\ref{sec:candidate} constructs the fixed self-conjugate partition
for each large degree. Two complementary sufficient tests show that
its square contains a prescribed target.

The first test, developed in Section~\ref{sec:band}, starts with one
tensor obtained by alternating along the candidate's rows in the first
factor and its columns in the second. A coordinate projection of its
diagonal orbit separates a smaller triangle from a width-two band and
two attached diagrams. The projected module is the full induction of
the triangle and complement modules. The cyclic staircase input supplies
every constituent on the triangle, so branching reduces the task to
finding a supported complement diagram contained in the desired target.
The band is an odd path of alternating pairs: its contraction is a word
of two-by-two matrices and adjugates. The attachments are handled by
nondegenerate symmetric forms and a parity-sensitive
Littlewood--Richardson splitting. Keeping the attached dual columns
away from the path endpoints makes every surviving term factor, so the
nonzero contractions do not cancel.

The second test, in Section~\ref{sec:balance}, works in the full square,
using column alternations in both factors. It allows independent
position permutations in the two layers before applying a common
projection. Label the two alphabets by positive integers; a pair of
letters $a,a'$ has output $a+a'-1$. Components can then be separated
by both the numbers and sums of their outputs.
This combines four-row path supports, two-row
symmetric-form supports, and an additional three-row support
construction. This
size-and-sum projection, as well as the explicit attachment criterion,
can be applied independently of the final degree estimates.

These tests divide the targets by their initial row and column sums.
The balance test supplies a target whenever its first four rows or first
four columns have small total size. For a remaining target, failure of
the band test in both orientations bounds the first few row sums and
column sums simultaneously. Section~\ref{sec:capacity} shows that, for
large parameters, a diagram satisfying both bounds has less area than
the target's prescribed size. Two exact capacity calculations settle
the remaining bounded parameters. The Murnaghan--Nakayama recursion of
Section~\ref{sec:finite}
verifies every eligible degree at most $64$. The latter computes exact
residues of whole Kronecker-coefficient vectors, not floating-point
approximations. We prove the soundness of every test and pruning rule
used by the programs and include their complete source. Thus the finite
computations close specified finite cases rather than impose an
unproved threshold on positivity.

Section~\ref{sec:prelim} fixes the common polytabloid conventions and the
exact cyclic companion input before either construction begins.
Section~\ref{sec:completion} assembles the candidate, the support tests,
and the finite range. Appendices~\ref{app:capacity-code}
and~\ref{app:smallcode} supply the two proof-critical programs; the
accompanying source also contains their independent checking tools.
